\documentclass[11pt,a4paper]{article}
\usepackage[utf8]{inputenc}
\usepackage[T1]{fontenc}
\usepackage[polish]{babel}
\usepackage{amsmath}
\usepackage{amsfonts}
\usepackage{amssymb}
\usepackage{graphicx}
\usepackage{booktabs}
\usepackage{listings}
\usepackage{xcolor}
\usepackage{hyperref}
\usepackage{float}
\usepackage{array}
\usepackage{multirow}
\usepackage{subcaption}
\usepackage{pdfpages}
\usepackage[a4paper,margin=2cm]{geometry}
\usepackage{booktabs}
\usepackage{longtable}
\usepackage{multirow}
\usepackage{needspace}
\usepackage{tikz}
\usepackage{algorithm}
\usepackage{algpseudocode}
\usepackage{pgfplots}  % Add this package for the axis environment
\pgfplotsset{compat=1.18}  % Set compatibility mode
\usetikzlibrary{shapes.geometric, arrows}

\lstset{
    basicstyle=\ttfamily\small,
    commentstyle=\color{gray},
    keywordstyle=\color{blue},
    stringstyle=\color{purple},
    numbers=left,
    numberstyle=\tiny\color{gray},
    breaklines=true,
    frame=single,
    captionpos=b
}

\title{Projekt Końcowy: Zmodyfikowany Algorytm Lasu Losowego z Adaptacyjnym Ważeniem Próbek}
\author{Jakub Szubzda - 331145 \\ Aleksandra Wasilewska - 331192}
\date{\today}

\begin{document}

\maketitle
\tableofcontents
\newpage

\section{Streszczenie założeń z Projektu Wstępnego}

Projekt dotyczył implementacji zmodyfikowanej wersji algorytmu generowania lasu losowego (Random Forest), w której do generowania kolejnych drzew częściej losowane są elementy ze zbioru uczącego, na których dotychczasowy model się mylił. Takie podejście łączy zalety standardowego lasu losowego z mechanizmem adaptacyjnego ważenia próbek znanym z algorytmów boostingowych jak AdaBoost.

Główne założenia projektu obejmowały:

\begin{itemize}
    \item Implementację standardowego algorytmu drzew decyzyjnych i lasu losowego od podstaw
    \item Wprowadzenie modyfikacji polegającej na adaptacyjnym ważeniu próbek błędnie klasyfikowanych
    \item Dodanie wariacji algorytmu pozwalających na wybór różnych strategii podziału, przycinania, głosowania i wyboru cech
    \item Przeprowadzenie eksperymentów na różnorodnych zbiorach danych
    \item Porównanie wydajności zmodyfikowanego algorytmu ze standardowym lasem losowym i AdaBoostu
\end{itemize}

Pierwotnie zaproponowano szereg parametrów konfiguracyjnych, takich jak:
\begin{itemize}
    \item \texttt{max\_features} - liczba cech branych do każdego drzewa
    \item \texttt{sample\_fraction} - część przykładów brana do każdego drzewa
    \item \texttt{n\_trees} - liczba drzew w lesie
    \item \texttt{max\_depth} - maksymalna głębokość drzewa
    \item \texttt{prune} - czy drzewa powinny być przycinane po wygenerowaniu
    \item \texttt{criterion} - kryterium podziału (gini lub entropy)
    \item \texttt{weighted\_voting} - czy głosy drzew są ważone względem ich dokładności
    \item \texttt{error\_weight\_increase} - wzrost wagi błędnie sklasyfikowanych przykładów
    \item \texttt{weighted\_feature\_sampling} - czy cechy są losowane proporcjonalnie do ich information gain
\end{itemize}

\section{Pełen opis funkcjonalny}

\subsection{Architektura zaimplementowanego rozwiązania}

W ramach projektu zaimplementowano kompleksową bibliotekę do tworzenia i trenowania zmodyfikowanego lasu losowego, składającą się z kilku modułów:

\begin{itemize}
    \item \texttt{ModifiedRandomForest.py} - główna klasa implementująca zmodyfikowany las losowy
    \item \texttt{tree.py} - implementacja drzewa decyzyjnego
    \item \texttt{split.py} - funkcje do znajdowania optymalnych podziałów
    \item \texttt{sampling.py} - funkcje do losowania próbek i normalizacji wag
    \item \texttt{pruning.py} - funkcje do przycinania drzew
    \item \texttt{prediction.py} - funkcje do generowania i agregowania predykcji
    \item \texttt{evaluation.py} - funkcje do oceny jakości modelu
\end{itemize}

Architektura rozwiązania pozwala na:
\begin{itemize}
    \item Tworzenie i trenowanie lasów losowych o różnych konfiguracjach
    \item Dostosowywanie strategii próbkowania i ważenia przykładów
    \item Eksperymentowanie z różnymi kryteriami podziału i strategiami głosowania
    \item Serializację i deserializację wytrenowanych modeli
    \item Obliczanie ważności cech
    \item Obsługę różnych rodzajów zbiorów danych
\end{itemize}

\subsection{Szczegółowy opis funkcjonalności}

\subsubsection{Klasa ModifiedRandomForest}

Główna klasa \texttt{ModifiedRandomForest} implementuje zmodyfikowany algorytm lasu losowego i udostępnia następujące funkcjonalności:

\begin{itemize}
    \item \texttt{fit(X, y)} - trenuje model na danych wejściowych
    \item \texttt{predict(X)} - generuje predykcje dla nowych danych
    \item \texttt{predict\_proba(X)} - generuje prawdopodobieństwa przynależności do klas
    \item \texttt{save\_model(filename)} - zapisuje model do pliku
    \item \texttt{load\_model(filename)} - wczytuje model z pliku
    \item \texttt{print\_mapping()} - wyświetla mapowanie indeksów klas na oryginalne etykiety
\end{itemize}

Model ten współpracuje z różnymi typami danych, obsługuje klasy nienumeryczne i automatycznie konwertuje etykiety klas na reprezentację wewnętrzną.

\subsubsection{Drzewa decyzyjne}

Zaimplementowana klasa \texttt{DecisionTree} stanowi podstawowy blok budulcowy lasu losowego:

\begin{itemize}
    \item Wspiera dwa kryteria podziału: Gini impurity i Information Gain (entropy)
    \item Umożliwia rekurencyjne budowanie drzewa z ograniczeniem maksymalnej głębokości
    \item Implementuje logikę przycinania drzew po ich zbudowaniu
    \item Oblicza ważność cech podczas trenowania
    \item Optymalizuje pamięć poprzez wykorzystanie \_\_slots\_\_ w węzłach drzew
    \item Umożliwia konwersję drzewa do formatu JSON i odtworzenie go z tego formatu
\end{itemize}

\subsubsection{Próbkowanie i ważenie}

Kluczową modyfikacją w stosunku do standardowego lasu losowego jest adaptacyjne ważenie próbek:

\begin{itemize}
    \item Początkowo wszystkie próbki mają równe wagi
    \item Po każdej iteracji (dodaniu nowego drzewa), wagi próbek, które zostały błędnie sklasyfikowane przez aktualny las, są zwiększane
    \item Wagi są normalizowane po każdej aktualizacji
    \item Wagi wpływają na prawdopodobieństwo wylosowania próbek do trenowania kolejnego drzewa
    \item Parametr \texttt{error\_weight\_increase} kontroluje siłę zwiększania wag
\end{itemize}

\subsubsection{Głosowanie}

Zaimplementowano dwa tryby głosowania przy generowaniu predykcji:

\begin{itemize}
    \item \textbf{Głosowanie większościowe:} każde drzewo ma równy głos
    \item \textbf{Głosowanie ważone:} głos każdego drzewa jest ważony jego dokładnością na zbiorze OOB (out-of-bag)
\end{itemize}

\subsubsection{Wybór cech}

Model umożliwia dwa sposoby wyboru cech przy budowie węzłów drzewa:

\begin{itemize}
    \item \textbf{Równomierne losowanie:} standardowy wybór losowy cech
    \item \textbf{Ważone losowanie cech:} losowanie cech z prawdopodobieństwem proporcjonalnym do ich information gain względem klasy
\end{itemize}

\section{Precyzyjny opis algorytmów oraz zbiorów danych}

\subsection{Algorytm zmodyfikowanego lasu losowego}

\subsubsection{Faza trenowania}

\begin{algorithm}
\caption{FitModifiedRandomForest(X, y, n\_trees, max\_features, sample\_fraction, max\_depth, prune, criterion, error\_weight\_increase, weighted\_feature\_sampling)}
\begin{algorithmic}[1]
\State $n\_samples \gets \text{length}(X)$
\State $n\_features \gets \text{length}(X[0])$
\State $forest \gets []$
\State $oob\_accuracies \gets []$
\State $feature\_importances \gets \text{vector}(n\_features, 0.0)$
\State $sample\_weights \gets \text{vector}(n\_samples, 1.0 / n\_samples)$

\For{$i \gets 1$ to $n\_trees$}
    \State $sample\_size \gets \text{round}(sample\_fraction \cdot n\_samples)$
    \State $idx \gets \text{sample\_with\_replacement}(0..n\_samples-1, size=sample\_size, prob=sample\_weights)$

    \State $oob\_mask \gets \text{vector}(n\_samples, \text{True})$
    \State $oob\_mask[idx] \gets \text{False}$
    \State $X\_train\_tree \gets X[idx]$
    \State $y\_train\_tree \gets y[idx]$
    \State $X\_oob \gets X[oob\_mask]$
    \State $y\_oob \gets y[oob\_mask]$

    \If{$weighted\_feature\_sampling$}
        \State $ig\_scores \gets \text{calculate\_information\_gain}(X\_train\_tree, y\_train\_tree)$
        \State $features\_subset \gets \text{weighted\_random\_selection}(max\_features, weights=ig\_scores)$
    \Else
        \State $features\_subset \gets \text{random\_selection}(n\_features, max\_features)$
    \EndIf

    \State $tree \gets \text{build\_tree}(X\_train\_tree, y\_train\_tree, features\_subset, max\_depth, criterion, min\_samples\_split, min\_samples\_leaf)$

    \If{$prune$}
        \If{$\text{length}(X\_oob) > 0$}
            \State $prune\_data\_X, prune\_data\_y \gets X\_oob, y\_oob$
        \Else
            \State $prune\_data\_X, prune\_data\_y \gets X\_train\_tree, y\_train\_tree$
        \EndIf
        \State $tree \gets \text{prune\_tree}(tree, prune\_data\_X, prune\_data\_y)$
    \EndIf

    \State $oob\_acc \gets 0.0$
    \If{$\text{length}(X\_oob) > 0$}
        \State $oob\_predictions \gets tree.predict(X\_oob)$
        \State $oob\_acc \gets \text{accuracy}(y\_oob, oob\_predictions)$
    \EndIf

    \State $tree\_importances \gets tree.feature\_importances\_$

    \State $forest.append(tree)$
    \State $oob\_accuracies.append(oob\_acc)$
    \State $feature\_importances \gets feature\_importances + tree\_importances$

    \If{$error\_weight\_increase > 0$}
        \State $current\_ensemble\_pred \gets \text{ensemble\_predict}(forest, X)$
        \State $errors \gets (current\_ensemble\_pred \neq y)$

        \For{$j \gets 0$ to $n\_samples-1$}
            \If{$errors[j]$}
                \State $sample\_weights[j] \gets sample\_weights[j] \cdot (1.0 + error\_weight\_increase)$
            \EndIf
        \EndFor

        \State $sample\_weights \gets \text{normalize}(sample\_weights)$
    \EndIf
\EndFor

\If{$\text{sum}(feature\_importances) > 0$}
    \State $feature\_importances \gets \text{normalize}(feature\_importances)$
\EndIf

\State \Return $forest, oob\_accuracies, feature\_importances$
\end{algorithmic}
\end{algorithm}

\subsubsection{Faza predykcji}

\begin{algorithm}
\caption{PredictModifiedRandomForest(forest, X, use\_weighted\_vote, oob\_accuracies)}
\begin{algorithmic}[1]
\State $all\_trees\_predictions \gets []$
\For{each $tree$ in $forest$}
    \State $tree\_predictions \gets tree.predict(X)$
    \State $all\_trees\_predictions.append(tree\_predictions)$
\EndFor

\If{$use\_weighted\_vote$}
    \State $final\_predictions \gets \text{perform\_weighted\_voting}(all\_trees\_predictions, oob\_accuracies)$
\Else
    \State $final\_predictions \gets \text{perform\_majority\_voting}(all\_trees\_predictions)$
\EndIf

\State \Return $final\_predictions$
\end{algorithmic}
\end{algorithm}

\subsubsection{Budowa drzewa decyzyjnego}

\begin{algorithm}
\caption{BuildTree(X, y, feature\_indices, depth, max\_depth, criterion, min\_samples\_split, min\_samples\_leaf)}
\begin{algorithmic}[1]
\State $n\_samples, \_ \gets X.shape$
\State $n\_classes \gets \text{length}(\text{unique}(y))$

\If{$\text{check\_stopping\_criteria}(depth, max\_depth, n\_samples, min\_samples\_split, n\_classes)$ \textbf{or} $\text{is\_pure}(y)$}
    \State \Return $\text{create\_leaf}(y)$
\EndIf

\State $feature\_idx, threshold, info\_gain, left\_indices, right\_indices \gets \text{find\_best\_split}(X, y, feature\_indices, criterion, min\_samples\_leaf)$

\If{$feature\_idx$ is None \textbf{or} $info\_gain \leq 0.0$}
    \State \Return $\text{create\_leaf}(y)$
\EndIf

\State $left\_tree \gets \text{BuildTree}(X[left\_indices], y[left\_indices], feature\_indices, depth+1, max\_depth, criterion, min\_samples\_split, min\_samples\_leaf)$

\State $right\_tree \gets \text{BuildTree}(X[right\_indices], y[right\_indices], feature\_indices, depth+1, max\_depth, criterion, min\_samples\_split, min\_samples\_leaf)$

\State \Return $\text{create\_node}(feature\_idx, threshold, left\_tree, right\_tree)$
\end{algorithmic}
\end{algorithm}

\subsection{Techniczne aspekty implementacji}

\subsubsection{Optymalizacja wydajności}

W celu przyspieszenia działania algorytmu, zastosowano szereg technik optymalizacji:

\begin{itemize}
    \item Wykorzystanie biblioteki Numba do kompilacji JIT funkcji krytycznych obliczeniowo
    \item Zastosowanie \_\_slots\_\_ w klasie Node dla efektywnego zarządzania pamięcią
    \item Wykorzystanie algorytmów przybliżonego podziału dla dużej liczby unikalnych wartości cechy
    \item Implementacja efektywnych metod normalizacji wag i ich aktualizacji
    \item Optymalizacja obliczania miar nieczystości (Gini, entropy) z użyciem funkcji wektoryzowanych
\end{itemize}

\subsubsection{Współpraca z różnymi typami danych}

Model został zaprojektowany z myślą o elastyczności i obsłudze różnych typów danych:

\begin{itemize}
    \item Automatyczna konwersja między różnymi formatami danych wejściowych
    \item Obsługa etykiet klas niebędących liczbami (stringi, kategorie)
    \item Serializacja i deserializacja z zachowaniem mapowania klas
    \item Wykrywanie typów cech i odpowiednia ich obsługa
\end{itemize}

\subsection{Opis zbiorów danych}

W eksperymentach wykorzystano 5 różnych zbiorów danych, różniących się wielkością, złożonością, typem cech i liczbą klas:

\begin{table}[H]
\caption{Charakterystyka zbiorów danych użytych w eksperymentach}
\label{tab:dataset_comparison}
\begin{tabular}{lccccc}
\toprule
Dataset & Próbki & Cechy & Klasy & Współczynnik niezbalansowania & Braki danych(\%) \\
\midrule
Iris & 150 & 4 & 3 & 1.0 & 0 \\
Wine Quality & 6497 & 11 & 7 & 567.2 & 0 \\
Adult Census & 48842 & 14 & 2 & 5.23 & 0.32 \\
Bank Marketing & 45211 & 16 & 2 & 7.55 & 7.21 \\
MNIST & 70000 & 784 & 10 & 1.25 & 0 \\
\bottomrule
\end{tabular}
\end{table}

\subsubsection{Iris Dataset}
Klasyczny zbiór danych zawierający pomiary trzech gatunków irysów. Cechy obejmują długość i szerokość działek kielicha oraz płatków korony. Jest to mały, zbalansowany zbiór danych z trzema klasami, często używany jako benchmark w uczeniu maszynowym.

\subsubsection{Wine Quality Dataset}
Zbiór danych zawierający informacje o chemicznych właściwościach win portugalskich (Vinho Verde) i ich jakości w skali od 3 do 9. Cechy obejmują między innymi kwasowość, zawartość alkoholu, cukru i kwasu cytrynowego. Jest to zbiór średniej wielkości z niezbalansowanymi klasami.

\subsubsection{Adult Census Dataset}
Zbiór danych pochodzący ze spisu ludności USA, zawierający informacje demograficzne i socjoekonomiczne. Zadanie polega na przewidywaniu, czy dochód osoby przekracza 50 000 dolarów rocznie. Cechy obejmują wiek, wykształcenie, zawód, stan cywilny i inne. Jest to duży zbiór danych z cechami mieszanymi (numeryczne i kategoryczne).

\subsubsection{Bank Marketing Dataset}
Zbiór danych dotyczący kampanii marketingowych banku, gdzie celem jest przewidywanie, czy klient zapisze się na lokatę terminową. Cechy obejmują dane demograficzne klientów, informacje o historii kontaktów i wskaźniki ekonomiczne. Jest to duży zbiór danych z silnym niezbalansowaniem klas.

\subsubsection{MNIST Dataset}
Zbiór danych zawierający obrazy odręcznie napisanych cyfr (0-9). Każdy obraz ma rozmiar 28x28 pikseli, co daje 784 cechy. Jest to duży zbiór danych z wieloma cechami i 10 klasami, stanowiący wyzwanie dla algorytmów nieneuronowych.

\section{Raport z przeprowadzonych testów oraz wnioski}

\subsection{Metodyka testowania}

Przeprowadzone eksperymenty podzielono na dwie główne fazy:
\begin{enumerate}
    \item \textbf{Faza 1:} Poszukiwanie optymalnych parametrów zmodyfikowanego algorytmu Random Forest dla wybranego zbioru danych (Wine Quality). Testowano wpływ poszczególnych parametrów na jakość modelu.

    \item \textbf{Faza 2:} Porównanie wydajności zmodyfikowanego Random Forest z algorytmami referencyjnymi (scikit-learn Random Forest i AdaBoost) na różnych zbiorach danych.
\end{enumerate}

Do oceny modeli wykorzystano następujące metryki:
\begin{itemize}
    \item Accuracy - ogólna dokładność klasyfikacji
    \item Precision - precyzja klasyfikacji
    \item Recall - czułość klasyfikacji
    \item F1 Score - średnia harmoniczna precision i recall
    \item ROC AUC - pole pod krzywą ROC
    \item Fit Time - czas trenowania modelu
\end{itemize}

\subsection{Wyniki eksperymentów fazy 1: Optymalizacja parametrów}

W pierwszej fazie eksperymentów zbadano wpływ różnych parametrów na wydajność modelu zmodyfikowanego Random Forest dla zbioru danych Wine Quality.

\subsubsection{Wpływ parametru max\_features}

\begin{table}[H]
\centering
\caption{Wpływ parametru max\_features na wydajność modelu}
\label{tab:max_features}
\begin{tabular}{lcccc}
\toprule
max\_features & Dokładność & F1 Score & ROC AUC & Czas trenowania (s) \\
\midrule
sqrt (3) & 0.617 & 0.603 & 0.776 & 43.52 \\
0.3 (3) & 0.617 & 0.603 & 0.776 & 37.23 \\
0.5 (5) & 0.617 & 0.607 & 0.773 & 54.85 \\
1.0 (11) & \textbf{0.639} & \textbf{0.630} & \textbf{0.783} & 110.20 \\
\bottomrule
\end{tabular}
\end{table}

Obserwacja: Wykorzystanie wszystkich cech (max\_features=1.0) dało najlepsze wyniki dla zbioru Wine Quality, choć kosztem znacznie dłuższego czasu trenowania. Wskazuje to na istotność wszystkich cech w tym zbiorze danych.

\subsubsection{Wpływ liczby drzew (n\_trees)}

\begin{table}[H]
\centering
\caption{Wpływ liczby drzew na wydajność modelu}
\label{tab:n_trees}
\begin{tabular}{lcccc}
\toprule
n\_trees & Dokładność & F1 Score & ROC AUC & Czas trenowania (s) \\
\midrule
10 & 0.569 & 0.554 & 0.740 & 6.28 \\
50 & 0.617 & 0.603 & 0.776 & 37.01 \\
100 & 0.634 & 0.620 & 0.774 & 95.47 \\
200 & \textbf{0.642} & \textbf{0.631} & 0.769 & 282.90 \\
\bottomrule
\end{tabular}
\end{table}

Obserwacja: Zwiększenie liczby drzew poprawia dokładność i F1 Score, ale z malejącym zyskiem. 100-200 drzew oferuje dobry kompromis między wydajnością a czasem obliczeniowym.

\subsubsection{Wpływ maksymalnej głębokości drzewa (max\_depth)}

\begin{table}[H]
\centering
\caption{Wpływ maksymalnej głębokości drzewa na wydajność modelu}
\label{tab:max_depth}
\begin{tabular}{lcccc}
\toprule
max\_depth & Dokładność & F1 Score & ROC AUC & Czas trenowania (s) \\
\midrule
5 & 0.411 & 0.409 & 0.635 & 16.24 \\
10 & 0.617 & 0.603 & 0.776 & 37.16 \\
None & \textbf{0.660} & \textbf{0.644} & \textbf{0.831} & 68.51 \\
\bottomrule
\end{tabular}
\end{table}

Obserwacja: Nieograniczona głębokość drzewa (max\_depth=None) dała najlepsze wyniki dla tego zbioru danych. Sugeruje to, że dla Wine Quality potrzebne są bardziej złożone drzewa do uchwycenia zależności.

\subsubsection{Wpływ zwiększania wagi błędnych przykładów (error\_weight\_increase)}

\begin{table}[H]
\centering
\caption{Wpływ parametru error\_weight\_increase na wydajność modelu}
\label{tab:error_weight_increase}
\begin{tabular}{lcccc}
\toprule
error\_weight\_increase & Dokładność & F1 Score & ROC AUC & Czas trenowania (s) \\
\midrule
0.0 & 0.587 & 0.538 & 0.782 & 27.03 \\
0.1 & \textbf{0.631} & \textbf{0.620} & \textbf{0.790} & 39.71 \\
0.3 & 0.617 & 0.603 & 0.776 & 36.96 \\
0.5 & 0.605 & 0.592 & 0.760 & 34.28 \\
1.0 & 0.589 & 0.573 & 0.729 & 30.82 \\
\bottomrule
\end{tabular}
\end{table}

Obserwacja: Umiarkowana wartość parametru error\_weight\_increase (0.1) dała najlepsze wyniki. Zbyt duże zwiększanie wag błędnie sklasyfikowanych przykładów może prowadzić do przeuczenia.

Na podstawie eksperymentów z fazy 1 wybrano następujące optymalne parametry do fazy 2:
\begin{itemize}
    \item max\_features = 1.0
    \item n\_trees = 200
    \item max\_depth = None
    \item error\_weight\_increase = 0.1
    \item weighted\_voting = True
\end{itemize}

Potwierdza to printout z fazy 2:

\begin{table}[H]
\centering
\caption{Wynik fazy 2 dla optymalnych parametrów na zbiorze Wine Quality}
\label{tab:phase2}
\begin{tabular}{lc}
\toprule
Metryka & Wartość \\
\midrule
Dokładność & 0.6854 \\
Precision & 0.6854 \\
Recall & 0.6854 \\
F1 Score & 0.6751 \\
ROC AUC & 0.8560 \\
Czas trenowania (s) & 1104.18 \\
\bottomrule
\end{tabular}
\end{table}

\subsection{Wyniki eksperymentów fazy 3: Porównanie modeli na różnych zbiorach danych}

W fazie 3 porównano zaimplementowany zmodyfikowany Random Forest (ModifiedRF) z algorytmami referencyjnymi: scikit-learn Random Forest (ScikitRF) i AdaBoost na pięciu zbiorach danych.

\subsubsection{Porównanie dokładności}

\begin{table}[H]
\centering
\caption{Porównanie dokładności (accuracy) modeli na różnych zbiorach danych}
\label{tab:accuracy_comparison}
\begin{tabular}{lccc}
\toprule
Dataset & ModifiedRF & ScikitRF & AdaBoost \\
\midrule
Iris & 0.900 & 0.900 & \textbf{0.933} \\
Wine Quality & 0.656 & \textbf{0.689} & 0.452 \\
Adult Census & 0.748 & \textbf{0.848} & 0.849 \\
Bank Marketing & \textbf{0.891} & 0.880 & 0.893 \\
MNIST & 0.878 & \textbf{0.921} & 0.628 \\
\bottomrule
\end{tabular}
\end{table}

\subsubsection{Porównanie F1 Score}

\begin{table}[H]
\centering
\caption{Porównanie F1 Score modeli na różnych zbiorach danych}
\label{tab:f1_comparison}
\begin{tabular}{lccc}
\toprule
Dataset & ModifiedRF & ScikitRF & AdaBoost \\
\midrule
Iris & 0.900 & 0.900 & \textbf{0.933} \\
Wine Quality & 0.640 & \textbf{0.677} & 0.398 \\
Adult Census & 0.764 & \textbf{0.845} & 0.840 \\
Bank Marketing & 0.857 & \textbf{0.858} & 0.863 \\
MNIST & 0.878 & \textbf{0.921} & 0.632 \\
\bottomrule
\end{tabular}
\end{table}

\subsubsection{Porównanie ROC AUC}

\begin{table}[H]
\centering
\caption{Porównanie ROC AUC modeli na różnych zbiorach danych}
\label{tab:roc_auc_comparison}
\begin{tabular}{lccc}
\toprule
Dataset & ModifiedRF & ScikitRF & AdaBoost \\
\midrule
Iris & 0.985 & 0.987 & 0.975 \\
Wine Quality & 0.831 & \textbf{0.859} & 0.630 \\
Adult Census & 0.861 & 0.895 & \textbf{0.902} \\
Bank Marketing & \textbf{0.712} & 0.678 & 0.738 \\
MNIST & 0.986 & \textbf{0.995} & 0.922 \\
\bottomrule
\end{tabular}
\end{table}

\subsubsection{Porównanie czasu trenowania}

\begin{table}[H]
\centering
\caption{Porównanie czasu trenowania (s) modeli na różnych zbiorach danych}
\label{tab:fit_time_comparison}
\begin{tabular}{lccc}
\toprule
Dataset & ModifiedRF & ScikitRF & AdaBoost \\
\midrule
Iris & 6.92 & \textbf{0.06} & 0.07 \\
Wine Quality & 75.34 & \textbf{1.02} & 0.25 \\
Adult Census & 137.65 & \textbf{2.21} & 2.39 \\
Bank Marketing & 122.88 & \textbf{1.73} & 1.16 \\
MNIST & 2052.59 & \textbf{22.61} & 25.29 \\
\bottomrule
\end{tabular}
\end{table}

\begin{figure}[H]
\centering
\begin{tikzpicture}
\begin{axis}[
    ybar,
    width=12cm,
    height=8cm,
    legend style={at={(0.5,1.05)}, anchor=south,legend columns=-1},
    ylabel={Accuracy},
    xlabel={Dataset},
    symbolic x coords={Iris,Wine,Adult,Bank,MNIST},
    xtick=data,
    nodes near coords,
    nodes near coords align={vertical},
    ymin=0.4, ymax=1.0,
    bar width=0.2cm,
]
\addplot coordinates {(Iris,0.900) (Wine,0.656) (Adult,0.748) (Bank,0.891) (MNIST,0.878)};
\addplot coordinates {(Iris,0.900) (Wine,0.689) (Adult,0.848) (Bank,0.880) (MNIST,0.921)};
\addplot coordinates {(Iris,0.933) (Wine,0.452) (Adult,0.849) (Bank,0.893) (MNIST,0.628)};
\legend{ModifiedRF, ScikitRF, AdaBoost}
\end{axis}
\end{tikzpicture}
\caption{Porównanie dokładności modeli na różnych zbiorach danych}
\label{fig:accuracy_comparison}
\end{figure}

\subsection{Wnioski z eksperymentów}

Na podstawie przeprowadzonych eksperymentów można sformułować następujące wnioski:

\begin{enumerate}
    \item \textbf{Wydajność modeli na różnych zbiorach danych:}
    \begin{itemize}
        \item Zmodyfikowany Random Forest (ModifiedRF) pokazał porównywalną wydajność do ScikitRF na zbiorze Bank Marketing, gdzie osiągnął najlepszą dokładność.
        \item Dla zbioru Iris wszystkie algorytmy osiągnęły wysoką dokładność, z najlepszym wynikiem dla AdaBoost.
        \item Na zbiorze MNIST, ModifiedRF uzyskał dobrą dokładność (87.8\%), choć niższą niż ScikitRF (92.1\%).
        \item AdaBoost radził sobie słabo na dużych zbiorach danych z wieloma cechami (np. MNIST) i na zbiorze Wine Quality.
    \end{itemize}

    \item \textbf{Wpływ parametrów na wydajność ModifiedRF:}
    \begin{itemize}
        \item Zwiększenie liczby drzew znacząco poprawia dokładność, choć z malejącymi korzyściami powyżej 100 drzew.
        \item Nieograniczona głębokość drzew (max\_depth=None) daje najlepsze wyniki dla złożonych zbiorów.
        \item Parametr error\_weight\_increase powinien być ustawiony umiarkowanie (najlepiej 0.1), aby uniknąć przeuczenia.
        \item Ważone głosowanie (weighted\_voting=True) poprawia wyniki w większości przypadków.
    \end{itemize}

    \item \textbf{Kompromisy wydajnościowe:}
    \begin{itemize}
        \item ModifiedRF wymaga znacznie więcej czasu obliczeniowego niż ScikitRF i AdaBoost (od 50 do 100 razy więcej).
        \item Zoptymalizowany ModifiedRF (z parametrami z fazy 2) poprawia dokładność o około 4 punkty procentowe w porównaniu do domyślnej konfiguracji, ale kosztem zwiększonych zasobów obliczeniowych.
    \end{itemize}

    \item \textbf{Zachowanie na różnych typach zbiorów:}
    \begin{itemize}
        \item ModifiedRF radzi sobie lepiej na zbiorach z niezbalansowanymi klasami (Bank Marketing) w porównaniu do ScikitRF.
        \item Na zbiorze z wieloma klasami i cechami (MNIST), ModifiedRF osiąga dobre wyniki, ale ustępuje wysoce zoptymalizowanej implementacji ScikitRF.
        \item AdaBoost jest najlepszy na małych, prostych zbiorach danych (Iris).
    \end{itemize}
\end{enumerate}

\section{Uwagi dotyczące wykonania projektu}

\subsection{Wyjaśnienie odchyleń od planu początkowego}

Podczas realizacji projektu napotkaliśmy na szereg wyzwań, które spowodowały, że musieliśmy wprowadzić pewne modyfikacje w stosunku do pierwotnych założeń projektu wstępnego:

\begin{enumerate}
    \item \textbf{Redukcja zakresu eksperymentów}: Z uwagi na wysoką złożoność obliczeniową zmodyfikowanego algorytmu Random Forest, nie byliśmy w stanie przeprowadzić wszystkich zaplanowanych kombinacji parametrów na wszystkich zbiorach danych. Zastosowaliśmy dwufazowe podejście, koncentrując się najpierw na jednym zbiorze (Wine Quality), a następnie testując optymalne parametry na pozostałych.

    \item \textbf{Priorytetyzacja wybranych parametrów}: Przeprowadzone eksperymenty pokazały, że niektóre parametry (np. \texttt{error\_weight\_increase}, \texttt{max\_features}, \texttt{max\_depth}) mają większy wpływ na wydajność modelu niż inne. Skupiliśmy się na dokładniejszym zbadaniu tych parametrów.

    \item \textbf{Zoptymalizowana implementacja}: Ze względu na ograniczenia wydajnościowe, wprowadziliśmy szereg optymalizacji kodu, w tym kompilację JIT z Numba dla krytycznych obliczeniowo funkcji oraz efektywne zarządzanie pamięcią z wykorzystaniem \_\_slots\_\_.

    \item \textbf{Modyfikacja zbiorów danych}: Dla zbioru MNIST zastosowaliśmy znaczącą redukcję wymiarowości, aby umożliwić efektywne trenowanie modelu na dostępnej infrastrukturze.

    \item \textbf{Ograniczenie liczby zbiorów danych}: W ramach projektu wstępnego przeprowadziliśmy szczegółową analizę większej liczby zbiorów danych niż te, które zostały ostatecznie wykorzystane w eksperymentach. Ze względu na czasochłonność trenowania modelów na wielu zbiorach, skupiliśmy się na wybranych 5 zbiorach danych, które umożliwiły wyciągnięcie najważniejszych wniosków. Szczegółowa analiza wszystkich rozważanych zbiorów danych znajduje się w dokumentacji wstępnej projektu.
\end{enumerate}

\subsection{Ograniczenia dotyczące zbiorów danych}

Warto podkreślić, że w fazie przygotowawczej przeanalizowano i przygotowano więcej zbiorów danych niż zostało ostatecznie użytych w eksperymentach. Ze względu na ograniczenia czasowe i obliczeniowe, szczególnie w kontekście trenowania modeli na dużych zbiorach z wieloma parametrami, zdecydowano się na skupienie na wybranych 5 zbiorach danych, które reprezentowały zróżnicowane scenariusze uczenia maszynowego:

\begin{itemize}
    \item Mały zbiór zbalansowany (Iris)
    \item Średniej wielkości zbiór z wieloma klasami (Wine Quality)
    \item Duże zbiory z cechami mieszanymi i niezbalansowaniem klas (Adult Census, Bank Marketing)
    \item Zbiór z dużą liczbą cech (MNIST)
\end{itemize}

Szczegółowe analizy statystyczne i wizualne wszystkich rozważanych zbiorów danych, w tym wykresy rozkładów, macierze korelacji, analizy brakujących wartości i miary ważności cech, zostały zawarte w dokumentacji projektu wstępnego. Przeprowadzone eksperymenty wykazały, że najcenniejsze wnioski dotyczące wydajności zmodyfikowanego algorytmu Random Forest można było uzyskać na podstawie wybranych zbiorów danych, więc rezygnacja z pozostałych nie wpłynęła znacząco na jakość badania.

\subsection{Szczegółowy opis przetwarzania danych}

Każdy z wykorzystanych zbiorów danych wymagał specyficznego przetwarzania wstępnego, dostosowanego do jego charakterystyk:

\begin{enumerate}
    \item \textbf{Iris Dataset}: Ze względu na niewielki rozmiar i brak brakujących wartości, zbiór ten wykorzystano w formie oryginalnej oraz w wersji z zastosowaniem adaptacyjnego kodowania percentylowego.

    \item \textbf{Wine Quality Dataset}: Zastosowano adaptacyjne kodowanie percentylowe dla cech numerycznych z liczbą przedziałów dostosowaną do zysku informacyjnego względem zmiennej celu.

    \item \textbf{Adult Census Dataset}: Dla tego zbioru usunięto cechę 'fnlwgt' jako nieistotną, oraz zastąpiono cechę 'education' przez 'education-num'. Utworzono też kilka wariantów:
    \begin{itemize}
        \item Podstawowy z usunięciem brakujących wartości i kodowaniem one-hot
        \item Z adaptacyjnym kodowaniem percentylowym cech numerycznych
        \item Z redukcją liczby kategorii (np. grupowanie kategorii workclass, education, native-country)
    \end{itemize}

    \item \textbf{Bank Marketing Dataset}: Usunięto cechy potencjalnie powodujące wyciek danych ('duration', 'pdays') oraz czasowe ('day\_of\_week', 'month', 'contact'). Utworzono nową cechę 'previously\_contacted' na podstawie 'previous'. Przygotowano warianty z imputacją brakujących wartości oraz z usunięciem niepełnych rekordów.
\end{enumerate}

\subsubsection{Szczegółowe przetwarzanie zbioru MNIST}

Ze względu na wysoki wymiar oryginalnego zbioru MNIST (784 cechy), zastosowano wielostopniową redukcję wymiarowości i inżynierię cech:

\begin{enumerate}
    \item \textbf{Filtracja wariancji}: Usunięto cechy o niskiej wariancji (poniżej 1\% maksymalnej wariancji), co wyeliminowało piksele brzegowe i niezmienne.

    \item \textbf{Usuwanie silnie skorelowanych cech}: Zastosowano analizę korelacji, eliminując cechy o współczynniku korelacji powyżej 0,95, co dodatkowo zmniejszyło liczbę pikseli.

    \item \textbf{Ekstrakcja cech z filtrami konwolucyjnymi}: Zastosowano 8 różnych filtrów konwolucyjnych:
    \begin{itemize}
        \item Filtry kierunkowe: poziomy, pionowy, dwa diagonalne
        \item Filtry wygładzające: uśredniający, gaussowski
        \item Filtry wyostrzające: wyostrzający, uwypuklający
    \end{itemize}
    Dla każdego filtra obliczono 5 statystyk: średnią, odchylenie standardowe, maksimum, 75. percentyl i liczbę wartości dodatnich, uzyskując 40 cech dla każdego obrazu.

    \item \textbf{Cechy regionowe}: Podzielono obrazy na 16 regionów (siatka 4×4) i obliczono dla każdego regionu 4 statystyki: średnią, odchylenie standardowe, maksimum i gęstość pikseli, otrzymując dodatkowe 64 cechy.

    \item \textbf{Cechy oparte na odległości}: Obliczono 8 cech opisujących właściwości geometryczne cyfry:
    \begin{itemize}
        \item Odległość środka masy od środka obrazu
        \item Rozkład masy w czterech kwadrantach obrazu
        \item Miary symetrii poziomej i pionowej
        \item Stosunek szerokości do wysokości
    \end{itemize}

    \item \textbf{Tworzenie wariantów zbioru danych}: Wygenerowano dwa warianty zbioru MNIST:
    \begin{itemize}
        \item Wariant pełny z kombinacją zredukowanych pikseli i cech inżynieryjnych
        \item Wariant zawierający tylko cechy inżynieryjne (112 cech), eliminujący całkowicie oryginalną reprezentację pikselową
    \end{itemize}
\end{enumerate}

Ta kompleksowa transformacja zbioru MNIST zredukowała liczbę cech z 784 do 112, zachowując istotne informacje dyskryminacyjne. Dzięki temu trenowanie lasu losowego stało się bardziej efektywne obliczeniowo, jednocześnie utrzymując wysoką dokładność klasyfikacji (87,8\% dla naszego modelu w porównaniu do 92,1\% dla zoptymalizowanej implementacji scikit-learn).

\subsection{Wpływ przetwarzania danych na wydajność modelu}

Nasze eksperymenty pokazały, że odpowiednie przetwarzanie danych miało istotny wpływ na wydajność modeli:

\begin{itemize}
    \item Adaptacyjne kodowanie percentylowe poprawiło odporność na wartości odstające i normalizację rozkładów
    \item Redukcja wymiarowości MNIST zmniejszyła czas trenowania z szacowanych kilkudziesięciu godzin do około 34 minut
    \item Imputacja brakujących wartości w zbiorze Bank Marketing zwiększyła dostępną liczbę próbek treningowych, poprawiając wydajność modelu
    \item Tworzenie nowych cech (jak 'previously\_contacted') zwiększyło moc predykcyjną modelu
\end{itemize}

Szczególnie warto podkreślić, że wygenerowane cechy dla zbioru MNIST wykazywały wyższy zysk informacyjny w stosunku do oryginalnych pikseli, co potwierdza skuteczność zastosowanych metod inżynierii cech. Metryki odległościowe i cechy oparte na filtrach konwolucyjnych okazały się szczególnie istotne dla rozróżniania cyfr.

\section{Podsumowanie}

Zrealizowany projekt obejmował implementację zmodyfikowanej wersji algorytmu Random Forest z adaptacyjnym ważeniem próbek, dostosowaną do różnych typów zbiorów danych klasyfikacyjnych. Zmodyfikowany algorytm wprowadza mechanizm zwiększania wag próbek błędnie klasyfikowanych przez las, co pozwala na lepsze dopasowanie do trudniejszych przypadków.

Główne osiągnięcia projektu obejmują:

\begin{itemize}
    \item Implementację pełnego algorytmu lasu losowego od podstaw, z wykorzystaniem technik optymalizacyjnych
    \item Dodanie mechanizmu adaptacyjnego ważenia próbek, inspirowanego algorytmami boostingowymi
    \item Opracowanie elastycznej biblioteki z konfigurowalnymi parametrami
    \item Przeprowadzenie obszernych eksperymentów na różnorodnych zbiorach danych
    \item Identyfikację optymalnych parametrów dla różnych scenariuszy
\end{itemize}

Przeprowadzone eksperymenty wykazały, że zmodyfikowany Random Forest może konkurować z istniejącymi implementacjami na wybranych zbiorach danych, szczególnie tych z niezbalansowanymi klasami. Jednak kosztem jest znacznie wyższe zapotrzebowanie na zasoby obliczeniowe. Dalsze prace mogłyby skupić się na optymalizacji wydajności algorytmu oraz eksploracji innych mechanizmów adaptacyjnego ważenia przykładów.

\end{document}
